\chapter{Consumo energetico}
\label{cap:consumi}

\begin{quote}\small
Questo capitolo risponde all'obiettivo O4. L'analisi è a titolo informativo e usa i
valori dichiarati dai produttori, perché non si dispone di strumentazione di misura
(né multimetro né resistenze di shunt), come concordato con il relatore. I consumi dei
tre sensori sono quelli riportati nelle schede tecniche del \cref{cap:moduli}. Qui si
aggiunge il microcontrollore, si sommano i contributi nelle configurazioni del nodo e
se ne ricavano l'autonomia a batteria e l'architettura che ne discende.
\end{quote}

\section{Dati di partenza}
\label{sec:consumi-dati}

Le correnti dei sensori sono già nelle specifiche del \cref{cap:moduli}
(\cref{tab:ld2410b-specs} per il \ldb, \cref{sec:ld2420} per il \ldc,
\cref{tab:pir-specs} per il PIR). Qui si riprendono soltanto i valori che entrano nei
conti. Il \ldb\ assorbe in media 80~mA a 5~V, circa 400~mW. Il \ldc\ assorbe 50~mA a
3,3~V, circa 165~mW. L'\pir\ assorbe meno di 50~\textmu A, circa 0,3~mW. Il rapporto
fra i due estremi è il dato strutturale di tutto il capitolo. Il radar consuma circa
1600 volte la corrente del PIR.

Il dimostratore del progetto monta un PIR digitale Panasonic
EKMB139~\cite{callisto2026dipme}, un componente della classe a bassissimo consumo per
i nodi a batteria, diverso dal modulo del kit. I valori di questo capitolo si
riferiscono al kit usato negli esperimenti. La conclusione sull'ordine di grandezza
del divario non dipende però da quale PIR si consideri.

Il microcontrollore è il componente che manca al \cref{cap:moduli}. La
\cref{tab:consumi-esp32} riporta le correnti dell'ESP32-WROOM-32 per modalità
operativa, dal datasheet Espressif~\cite{esp32datasheet}.

\begin{table}[htbp]
\centering
\caption{Consumo dell'ESP32-WROOM-32 per modalità operativa, dal datasheet
Espressif~\cite{esp32datasheet}.}
\label{tab:consumi-esp32}
\small
\begin{tabular}{@{}lr@{}}
\toprule
\textbf{Modalità} & \textbf{Corrente} \\
\midrule
Attivo, WiFi in trasmissione (picco) & 160--260~mA \\
Attivo, WiFi connesso in ascolto & 95--100~mA \\
Modem-sleep (CPU attiva a 80--240~MHz, WiFi spento) & 20--68~mA \\
Light-sleep & 0,8~mA \\
Deep-sleep (solo dominio RTC) & 0,01~mA \\
\bottomrule
\end{tabular}
\end{table}

\section{Il nodo completo}
\label{sec:consumi-nodi}

Sommando microcontrollore e sensore si ottengono le configurazioni usate in questa
tesi e quelle rilevanti per il progetto (\cref{tab:consumi-nodi}). Le prime due sono
il logger e la web UI dei \cref{cap:moduli,cap:webui}. Le ultime due sono le
configurazioni a cui il nodo del progetto può aspirare.

\begin{table}[htbp]
\centering
\caption{Consumo stimato del nodo completo nelle configurazioni di interesse, somma
dei valori di datasheet.}
\label{tab:consumi-nodi}
\small
\begin{tabular}{@{}lrr@{}}
\toprule
\textbf{Configurazione} & \textbf{Corrente} & \textbf{Potenza a 5~V} \\
\midrule
ESP32 attivo + \ldb\ + WiFi (web UI) & 180--260~mA & 0,9--1,3~W \\
ESP32 attivo + \ldb\ (logger seriale) & 110--150~mA & 0,55--0,75~W \\
ESP32 attivo + \ldc & 80--120~mA & 0,4--0,6~W \\
ESP32 attivo + PIR & 30--70~mA & 0,15--0,35~W \\
ESP32 in deep-sleep + PIR, sveglia su interrupt & $\mathbf{0{,}06}$~\textbf{mA} & 0,3~mW \\
\bottomrule
\end{tabular}
\end{table}

L'ultima riga è il punto chiave. L'uscita digitale del PIR può risvegliare l'ESP32 dal
deep-sleep con un interrupt su GPIO. Il nodo può quindi restare in attesa a costo
praticamente nullo, con il PIR nel ruolo di guardiano. Il radar costa tre ordini di
grandezza in più ed è progettato per un'alimentazione fissa.

\section{Autonomia a batteria}
\label{sec:autonomia}

Con una cella litio-ione 18650 da 3000~mAh e un rendimento del regolatore dell'85\,\%
si ottengono le autonomie di \cref{tab:autonomia}. La \cref{fig:consumi} le riassume
insieme ai consumi dei componenti.

\begin{table}[htbp]
\centering
\caption{Autonomia stimata con una cella Li-ion 18650 da 3000~mAh e rendimento del
regolatore dell'85\,\%. Nel regime dormiente il fattore limitante è l'autoscarica
della cella, non il consumo.}
\label{tab:autonomia}
\small
\begin{tabular}{@{}lrr@{}}
\toprule
\textbf{Scenario} & \textbf{Corrente media} & \textbf{Autonomia} \\
\midrule
Nodo mmWave attivo con WiFi (\ldb) & 200~mA & 13~h \\
Nodo mmWave attivo senza WiFi continuo & 130~mA & 20~h \\
Nodo PIR sempre attivo (ESP32 in modem-sleep) & 25~mA & 4 giorni \\
Nodo dormiente (deep-sleep, PIR come guardiano) & 0,06~mA & 4--5 anni \\
\bottomrule
\end{tabular}
\end{table}

\begin{figure*}[htbp]
\forceversofloat
\centering
\includegraphics[width=\linewidth]{figures/fig13_consumi}
\caption{Consumi da datasheet dei singoli componenti e autonomia stimata dei nodi
completi, in scala logaritmica. Il PIR può fare da sveglia a costo quasi nullo. Il
radar costa tre ordini di grandezza in più ed è pensato per alimentazione fissa.}
\label{fig:consumi}
\end{figure*}

\section{Implicazione per DIPME: l'architettura ibrida}
\label{sec:architettura-ibrida}

I numeri precedenti risolvono un'apparente contraddizione del progetto. Il nodo
integrato negli arredi deve restare in attesa per anni in tempo di pace e deve
funzionare a batteria durante il blackout che segue il sisma. I due requisiti sono
incompatibili con un sensore da 80~mA sempre acceso, ma anche con un sensore che non
vede le persone ferme (\cref{cap:confronto}). La soluzione non è scegliere fra i due
sensori. È metterli in sequenza temporale.

\begin{enumerate}[leftmargin=1.8em,itemsep=3pt]
  \item \textbf{Tempo di pace.} L'intero nodo in deep-sleep, con il PIR spento o usato
        come guardiano. Consumo dell'ordine dei microampere, autonomia pluriennale
        limitata dall'autoscarica della batteria.
  \item \textbf{Passaggio in modalità emergenza.} L'accelerometro sul gateway o il
        rilevamento del blackout risveglia l'ESP32. Il PIR può contribuire come sorgente
        di interrupt a costo nullo.
  \item \textbf{Emergenza.} Si accende il radar per il rilevamento fine (persona ferma
        sotto il banco, respiro, indice di vitalità). L'autonomia di 13--20~ore copre la
        finestra critica iniziale delle operazioni di ricerca e soccorso.
  \item \textbf{Estensione della finestra.} Per coprire le prime 72~ore, l'orizzonte
        convenzionale del soccorso, si può maggiorare la batteria oppure alternare il
        radar con un ciclo di lavoro. Un minuto acceso ogni cinque riduce il consumo
        medio di circa cinque volte e porta l'autonomia oltre le 72~ore. Il prezzo è
        una latenza di rilevamento fino a quattro minuti, accettabile in questo
        scenario, perché la domanda è se c'è qualcuno sotto il banco e non in quale
        istante si è mosso.
\end{enumerate}

Il PIR quindi non va sostituito. Non è un concorrente del radar ma il guardiano a
basso costo che decide quando accenderlo. Il radar è lo strumento di misura della fase
di emergenza. I dati del \cref{cap:confronto} e di questo capitolo lo mostrano da due
direzioni opposte. Il \cref{cap:conclusioni} riprende questa architettura fra le
raccomandazioni per il progetto.
